(1) 一个无界均匀、不可压、密度 \(\rho\) 的理想流体在角速度 \(\boldsymbol\Omega=\Omega\mathbf e_z\) 的坐标系中静止。忽略离心静力压强对扰动的贡献，令 \(\mathbf u,\pi\) 为微小扰动速度和压强，建立线性Euler及不可压条件，并以波矢 \(\mathbf k=k(\sin\theta,0,\cos\theta)\) 的平面波推导两支惯性波色散 \(\omega(\mathbf k)\)。
(2) 构造与 \(\mathbf k\) 正交的实正交基 \(\mathbf e_1=(\cos\theta,0,-\sin\theta)\)、\(\mathbf e_2=(0,1,0)\)。求正频率分支的准确复速度偏振和压力振幅，说明惯性波为何为横向圆极化而非纵声波；讨论 \(\theta=0,\pi/2\) 的特殊情况。
(3) 由色散关系独立计算相速度向量与群速度向量，证明所有非零频率惯性波的群速度严格垂直波矢，并给出群速大小。讨论较大波数时相速与群速的尺度律及局域线性模型适用条件。
(4) 从线性Euler方程推导局部动能守恒及压力功率流。对第二问的平面波明确算出周期平均能量密度与平均能流，证明二者之比恰等于第三问群速度，并说明旋转科里奥利力不会直接做功。
